\subsection{Optimization Proposal for High-Dimensional Challenges}

To address the inefficiency of standard Monte Carlo sampling in higher dimensions, this section evaluates two sampling-based optimization techniques: Quasi-Monte Carlo and Markov Chain Monte Carlo. These methods are compared under the same geometric setting used in the baseline Monte Carlo experiment, namely the balls $A_n$ and $B_n$.

The baseline Monte Carlo estimator is given by
\[
\widehat{V}_{MC}
=
V_{\text{box}}
\cdot
\frac{1}{N}
\sum_{i=1}^{N}
\mathbf{1}\{x_i \in \mathcal{B}_n\},
\]
where $V_{\text{box}}$ is the volume of the bounding box, $N$ is the number of sampled points, and $\mathbf{1}\{x_i \in \mathcal{B}_n\}$ is an indicator function equal to $1$ if the point lies inside the target ball and $0$ otherwise.

Although this estimator is simple and intuitive, its efficiency decreases in higher dimensions because the ratio between the volume of the ball and the volume of the bounding box becomes small. As a result, many sampled points do not contribute to the estimate.

\subsubsection{Quasi-Monte Carlo Method}

The first optimization technique is Quasi-Monte Carlo sampling using a Sobol low-discrepancy sequence. Unlike standard Monte Carlo, which uses pseudo-random samples, Quasi-Monte Carlo generates points that are more evenly distributed across the sampling domain.

The Quasi-Monte Carlo estimator has the same volume-ratio structure:
\[
\widehat{V}_{QMC}
=
V_{\text{box}}
\cdot
\frac{1}{N}
\sum_{i=1}^{N}
\mathbf{1}\{s_i \in \mathcal{B}_n\},
\]
where $s_i$ denotes the $i$-th Sobol point mapped from $[0,1]^n$ to the corresponding bounding box.

For $A_n$, the Sobol points are mapped to $[0,1]^n$, while for $B_n$, they are mapped to $[-1,1]^n$.

The main advantage of Quasi-Monte Carlo is that it reduces random clustering and improves coverage of the sampling space. This can lead to more stable estimates than ordinary Monte Carlo, especially in low and moderate dimensions.

\subsubsection{Markov Chain Monte Carlo Method}

The second optimization technique is Markov Chain Monte Carlo using a random-walk proposal. Instead of generating independent samples, Markov Chain Monte Carlo constructs a sequence of dependent samples:
\[
X_{t+1} = X_t + \epsilon_t,
\]
where $\epsilon_t$ is drawn from a proposal distribution.

In this experiment, a normal random-walk proposal is used:
\[
\epsilon_t \sim \mathcal{N}(0, \sigma^2 I).
\]

A proposed point is accepted if it remains inside the bounding box. Otherwise, the chain stays at the current point. Therefore, the Markov chain targets the uniform distribution over the bounding box.

The Markov MC volume estimator is
\[
\widehat{V}_{MMC}
=
V_{\text{box}}
\cdot
\frac{1}{T}
\sum_{t=1}^{T}
\mathbf{1}\{X_t \in \mathcal{B}_n\},
\]
where $T$ is the number of retained samples after burn-in.

Unlike standard Monte Carlo and Quasi-Monte Carlo, the samples generated by Markov MC are correlated. Therefore, Markov MC is not always expected to outperform independent sampling methods. However, it provides an alternative sampling strategy and allows us to study how proposal tuning affects sampling efficiency.

\subsubsection{Proposal Tuning for Markov MC}

For the Markov MC method, proposal tuning is important because the proposal scale controls the movement of the Markov chain. If the scale is too small, the chain moves slowly and produces highly correlated samples. If the scale is too large, many proposed moves are rejected.

To evaluate the proposal quality, we compare different proposal distributions and scales using two metrics: acceptance rate and effective sample size (ESS). The acceptance rate measures how often proposed moves are accepted, while ESS measures the amount of useful independent information contained in the correlated Markov chain samples.

The tuning results show that the normal proposal with a moderate scale provides a reasonable balance between sample movement and acceptance rate. Although a higher acceptance rate may appear desirable, it does not always imply better performance because very small proposal steps can cause slow exploration. Therefore, the selected proposal should balance stability and movement.

\subsubsection{Numerical Comparison}

The following table summarizes the numerical results for $A_2$, $B_2$, $A_3$, and $B_3$. The estimated volumes are compared with the exact theoretical volumes.

\begin{table}[h]
\centering
\caption{Comparison of optimization methods for volume estimation.}
\label{tab:optimization_results}
\begin{tabular}{l l r r r r}
\hline
Shape & Method & Estimated Volume & Exact Volume & Absolute Error & Relative Error (\%) \\
\hline
$A_2$ & Quasi-Monte Carlo & -- & -- & -- & -- \\
$A_2$ & Markov MC & -- & -- & -- & -- \\
$B_2$ & Quasi-Monte Carlo & -- & -- & -- & -- \\
$B_2$ & Markov MC & -- & -- & -- & -- \\
$A_3$ & Quasi-Monte Carlo & -- & -- & -- & -- \\
$A_3$ & Markov MC & -- & -- & -- & -- \\
$B_3$ & Quasi-Monte Carlo & -- & -- & -- & -- \\
$B_3$ & Markov MC & -- & -- & -- & -- \\
\hline
\end{tabular}
\end{table}

Overall, Quasi-Monte Carlo improves the spatial coverage of the sampled points and usually gives more stable estimates than pseudo-random Monte Carlo in low dimensions. Markov MC provides a useful alternative sampling mechanism, but its samples are correlated, so its performance depends strongly on proposal tuning. Therefore, Quasi-Monte Carlo is the more direct optimization for this volume estimation problem, while Markov MC is useful for demonstrating how Markov-chain-based sampling can be applied and evaluated.
